\begin{center}
	\large{\textbf{KATA PENGANTAR}}
\end{center}
\indent Alhamdulillahi Rabbil 'Alamin. Segala puja dan puji syukur dipanjatkan kehadirat Allah SWT karena telah melimpahkan rahmat, hidayah, petunjuk, kekuatan, dan rezeki kepada penulis. Tak lupa sholawat serta salam penulis haturkan kepada junjungan dan panutan mulia Nabi Muhammad SAW yang membawa dunia dari jalan kegelapan menuju jalan yang terang benderang. Dengan demikian, berkat rahmat Allah SWT serta risalah Nabi Muhammad SAW, penulis dapat menyelesaikan Tugas Akhir yang berjudul:
\begin{center}
	\section*{OPTIMASI PORTOFOLIO BERBASIS MODEL MARKOWITZ MENGGUNAKAN \textit{VARIATIONAL QUANTUM EIGENSOLVER}}
\end{center}
sebagai salah satu syarat kelulusan Program Sarjana Fisika FSAD ITS.\\
\indent Penulisan Tugas Akhir ini dapat terselesaikan dengan baik berkat hadirnya bantuan dan bimbingan dari berbagai pihak. Penulis menyampaikan terima kasih kepada semua pihak yang telah membantu atas terselesaikannya Tugas Akhir ini.
\begin{enumerate}
    \item Tuhan yang maha Esa karena tanpa kasih sayang dan rahmat-Nya, penulis tidak akan bisa sampai di titik ini.
    \item Ibu dan Ayah yang senantiasa memberikan dukungan, doa, dan harapan kepada penulis.
    \item Bapak Dr. Heru Sukamto, M.Si. selaku dosen pembimbing I dan Bapak Dr. Gagus K. Sunnardianto selaku dosen pembimbing II yang telah meluangkan waktu, tenaga, dan pikiran untuk membimbing penulis.
	\item Bapak Prof. Agus Purwanto, D.Sc. dan Ibu Fahmi Astuti, M.Si., Ph.D. selaku dosen penguji yang telah meluangkan energi dan ilmunya sehingga penulis belajar untuk menjadi lebih teliti dan belajar untuk dapat menyampaikan ilmu agar lebih mudah dimengerti.
	\item Teman-teman yang meluangkan waktunya untuk hadir pada sidang Tugas Akhir mulai dari sidang terbuka hingga sidang "tertutup" tanggal 9 Juli 2026 karena doa, kehadiran, dan tawa para audiens saat sidang sangat membantu meringankan beban penulis dalam menjawab pertanyaan dari dosen penguji. 
    \item Teman-teman Fisika ITS angkatan 2022 \textit{Betelgeuse} yang menjadi keluarga penulis selama berkuliah di Fisika ITS. Terima kasih penulis khususnya kepada teman-teman kelompok belajar \textit{RUDAL.CORP} karena telah mewarnai perjuangan selama berkuliah dan membantu penulis dalam menjalani perkuliahan sehingga tidak ada mata kuliah yang mengulang sepanjang masa kuliah. 
	\item Anggota Laboratorium Fisika Teori dan Fislsafat Alam (LaFTiFA) sebagai tempat penulis berkembang dan belajar selama satu semester penuh. Terkhusus Ilham yang sudah menjerumuskan penulis ke LaFTiFA dan menjadi partner dari maba hingga tugas akhir, Alief yang sudah menjadi "dosen pembimbing III" yang selalu menuntut penulis untuk belajar hal-hal penting dalam menyusun tugas akhir, Lail yang telah meminjamkan catatan kuliah teorinya agar penulis dapat mengejar ketertinggalan mata kuliah teori, dan Agung yang telah menjadi teman diskusi dalam merencanakan masa depan.
	\item Seluruh pihak yang tidak bisa disebuthkan satu-persatu yang selalu memberikan dukungan, doa, dan harapan dalam menyelesaikan tugas akhir ini.
\end{enumerate}
Penulis menyadari bahwa dalam penulisan ini masih terdapat kekurangan, sehingga segala kritik dan saran yang sifatnya membangun sangat diharapkan untuk kesempurnaan Tugas Akhir ini. Penulis berharap  Tugas Akhir ini dapat bermanfaat bagi penulis sendiri pada khususnya dan pembaca pada umumnya.

\begin{flushright}
	Surabaya, \today \\
	\vspace{2cm}
	Penulis
\end{flushright}